\documentclass[11pt,a4paper]{article}

\usepackage{amsmath,amssymb}
\usepackage{graphicx}
\usepackage{booktabs}
\usepackage{hyperref}
\usepackage[margin=2.5cm]{geometry}
\usepackage{bm}
\usepackage{algorithm}
\usepackage{algpseudocode}

\title{Stereo Ensemble PIV for Reynolds Stress Extraction:\\Method and Validation}
\author{Technical Documentation}
\date{\today}

\begin{document}

\maketitle

\begin{abstract}
This document describes a method for extracting the Reynolds stress tensor from stereo particle image velocimetry (PIV) data using ensemble correlation analysis. The approach exploits the relationship between correlation peak widths and velocity fluctuation statistics: the auto-correlation width reflects particle image geometry, while the cross-correlation width includes additional broadening due to turbulent velocity fluctuations. By simultaneously fitting both correlations with a coupled 3D Gaussian model, we extract the Reynolds stress tensor $R_{ij} = \langle u'_i u'_j \rangle$. We validate the method using synthetic stereo images with known ground-truth Reynolds stresses, demonstrating recovery with approximately 30\% relative error for 2000 image pairs.
\end{abstract}

\section{Introduction}

Particle Image Velocimetry (PIV) is a cornerstone technique for measuring fluid velocity fields. While standard PIV cross-correlation provides instantaneous velocity measurements, extracting turbulence statistics such as the Reynolds stress tensor traditionally requires computing statistics over many independent realizations---a computationally expensive process that discards information contained in the correlation peak shape.

The \textit{ensemble correlation} approach offers an alternative: by averaging correlation functions across many image pairs before peak detection, information about velocity fluctuations is encoded directly in the correlation peak width. Specifically:
\begin{itemize}
    \item The \textbf{auto-correlation} peak width reflects only the particle image shape (geometric spread)
    \item The \textbf{cross-correlation} peak width reflects particle shape \textit{plus} turbulent broadening
\end{itemize}

The difference between these widths yields the Reynolds stress tensor. This document describes the complete pipeline from synthetic image generation through Reynolds stress extraction.

\section{Synthetic Data Generation}

\subsection{Stereo Camera Geometry}

We employ a symmetric stereo configuration with two cameras positioned at angles $\pm\theta$ from the optical axis (typically $\theta = 45°$). Each camera follows a pinhole projection model.

For a camera at position $\mathbf{p}_c$ with optical axis direction $\hat{\mathbf{n}}$, the projection of a 3D point $\mathbf{X} = (X, Y, Z)^T$ onto image coordinates $(u, v)$ is:
\begin{equation}
    \begin{pmatrix} u \\ v \end{pmatrix} = \frac{f}{Z'} \begin{pmatrix} X' \\ Y' \end{pmatrix} + \begin{pmatrix} c_x \\ c_y \end{pmatrix}
\end{equation}
where $(X', Y', Z')^T = \mathbf{R}(\mathbf{X} - \mathbf{p}_c)$ are coordinates in the camera frame, $f$ is the focal length, and $(c_x, c_y)$ is the principal point.

The camera rotation matrix $\mathbf{R}$ aligns the camera's local $z$-axis with the optical axis direction:
\begin{equation}
    \mathbf{R} = \begin{pmatrix}
        \cos\theta & 0 & \sin\theta \\
        0 & 1 & 0 \\
        -\sin\theta & 0 & \cos\theta
    \end{pmatrix}
\end{equation}
for a camera rotated by angle $\theta$ about the $y$-axis.

\subsection{Particle Field Generation}

Particles are distributed uniformly within a 3D measurement volume of size $L_x \times L_y \times L_z$. For frame A, particle positions are:
\begin{equation}
    \mathbf{x}_A^{(n)} \sim \mathcal{U}\left(-\frac{L}{2}, \frac{L}{2}\right)^3, \quad n = 1, \ldots, N_p
\end{equation}

Particle displacements between frames A and B are drawn from a multivariate normal distribution:
\begin{equation}
    \Delta\mathbf{x}^{(n)} = \mathbf{x}_B^{(n)} - \mathbf{x}_A^{(n)} \sim \mathcal{N}(\bm{\mu}, \mathbf{R})
\end{equation}
where $\bm{\mu}$ is the mean displacement (mean velocity times $\Delta t$) and $\mathbf{R}$ is the Reynolds stress tensor:
\begin{equation}
    \mathbf{R} = \begin{pmatrix}
        \langle u'u' \rangle & \langle u'v' \rangle & \langle u'w' \rangle \\
        \langle v'u' \rangle & \langle v'v' \rangle & \langle v'w' \rangle \\
        \langle w'u' \rangle & \langle w'v' \rangle & \langle w'w' \rangle
    \end{pmatrix}
\end{equation}

This formulation directly encodes the turbulence statistics we wish to recover.

\subsection{Image Rendering}

Each particle projects onto the camera sensor as a Gaussian intensity distribution:
\begin{equation}
    I(u, v) = I_0 \exp\left(-\frac{(u - u_0)^2 + (v - v_0)^2}{2\sigma_p^2}\right)
\end{equation}
where $(u_0, v_0)$ is the projected particle center and $\sigma_p$ controls the particle image diameter (typically 2--3 pixels for optimal PIV).

The total image is the sum of all particle contributions:
\begin{equation}
    I_{\text{total}}(u, v) = \sum_{n=1}^{N_p} I^{(n)}(u, v)
\end{equation}

\section{Volume Reconstruction: MinLOS}

\subsection{MLOS Principle}

To perform 3D correlation, we first reconstruct a volumetric intensity field from the stereo images using the Multiplicative Line-of-Sight (MLOS) approach. For each voxel $\mathbf{x}$ in the reconstruction volume:

\begin{enumerate}
    \item Project $\mathbf{x}$ to pixel coordinates $(u_c, v_c)$ in each camera $c$
    \item Sample the image intensity $I_c(u_c, v_c)$ using bilinear interpolation
    \item Combine intensities from all cameras
\end{enumerate}

The standard MLOS combination is multiplicative:
\begin{equation}
    V_{\text{MLOS}}(\mathbf{x}) = \prod_{c=1}^{N_c} I_c(u_c, v_c)
\end{equation}

Real particles appear bright in \textit{all} cameras, giving a high product. Empty space has low intensity in at least one view, giving a low product.

\subsection{Ghost Particles}

A fundamental limitation of tomographic reconstruction is \textit{ghost particles}: spurious high-intensity regions where lines of sight from different cameras intersect incorrectly. With only two cameras, ghost artifacts can be significant.

\subsection{MinLOS for Ghost Suppression}

We employ the \textit{Minimum} Line-of-Sight (MinLOS) variant to suppress ghosts:
\begin{equation}
    V_{\text{MinLOS}}(\mathbf{x}) = \min_{c} I_c(u_c, v_c)
\end{equation}

Ghost rays appear bright in one camera but correspond to empty space in other views. The minimum operation selects the low intensity, effectively suppressing ghosts while preserving real particles (which are bright in all views).

\section{Ensemble Correlation Method}

\subsection{3D FFT-Based Correlation}

The cross-correlation between volumes $A$ and $B$ is computed efficiently using the FFT:
\begin{equation}
    R_{AB}(\bm{\Delta}) = \mathcal{F}^{-1}\left[\mathcal{F}(A) \cdot \mathcal{F}(B)^*\right]
\end{equation}
where $\mathcal{F}$ denotes the 3D Fourier transform and $^*$ denotes complex conjugation.

Similarly, the auto-correlation is:
\begin{equation}
    R_{AA}(\bm{\Delta}) = \mathcal{F}^{-1}\left[|\mathcal{F}(A)|^2\right]
\end{equation}

\subsection{Ensemble Averaging in Frequency Domain}

For computational efficiency, we accumulate correlations in the frequency domain:
\begin{align}
    \sum_k \hat{R}_{AA}^{(k)} &= \sum_k \mathcal{F}(A_k) \cdot \mathcal{F}(A_k)^* \\
    \sum_k \hat{R}_{AB}^{(k)} &= \sum_k \mathcal{F}(A_k) \cdot \mathcal{F}(B_k)^*
\end{align}

The ensemble-averaged correlations are then:
\begin{equation}
    \langle R_{AA} \rangle = \mathcal{F}^{-1}\left[\frac{1}{N}\sum_k \hat{R}_{AA}^{(k)}\right], \quad
    \langle R_{AB} \rangle = \mathcal{F}^{-1}\left[\frac{1}{N}\sum_k \hat{R}_{AB}^{(k)}\right]
\end{equation}

\subsection{Background Subtraction}

Ghost particles and reconstruction artifacts create spurious correlations that persist across the ensemble. We remove these using background subtraction:
\begin{align}
    \tilde{R}_{AA} &= \langle A \star A \rangle - \langle A \rangle \star \langle A \rangle \\
    \tilde{R}_{AB} &= \langle A \star B \rangle - \langle A \rangle \star \langle B \rangle
\end{align}

The background term $\langle A \rangle \star \langle B \rangle$ (correlation of mean volumes) captures static artifacts that should be subtracted.

\subsection{Physical Interpretation of Correlation Widths}

The key insight enabling Reynolds stress extraction is the relationship between correlation peak shapes and underlying statistics:

\textbf{Auto-correlation} ($A \star A$): Correlating a volume with itself yields a peak whose width reflects only the \textit{geometric spread} of particle images---the point-spread function of the imaging system:
\begin{equation}
    \bm{\Sigma}_{\text{auto}} = \bm{\Sigma}_{\text{geo}}
\end{equation}

\textbf{Cross-correlation} ($A \star B$): Correlating frame A with displaced frame B yields a peak broadened by both geometry \textit{and} the distribution of particle displacements (turbulent fluctuations):
\begin{equation}
    \bm{\Sigma}_{\text{cross}} = \bm{\Sigma}_{\text{geo}} + \bm{\Sigma}_{\text{turb}}
\end{equation}

Therefore, the turbulent contribution (Reynolds stress) is:
\begin{equation}
    \boxed{\bm{\Sigma}_{\text{turb}} = \bm{\Sigma}_{\text{cross}} - \bm{\Sigma}_{\text{geo}} = \mathbf{R}}
\end{equation}

\section{Reynolds Stress Extraction via 3D Gaussian Fitting}

\subsection{Stacked Gaussian Model}

We fit a 3D Gaussian model simultaneously to both correlation volumes with coupled covariance constraints. The model has 22 parameters:

\begin{table}[h]
\centering
\begin{tabular}{@{}lll@{}}
\toprule
Parameters & Count & Description \\
\midrule
$A_{\text{auto}}, A_{\text{cross}}$ & 2 & Amplitudes \\
$B_{\text{auto}}, B_{\text{cross}}$ & 2 & Background offsets \\
$\bm{\Sigma}_{\text{geo}}$ & 6 & Geometric covariance (symmetric $3\times3$) \\
$\bm{\Sigma}_{\text{turb}}$ & 6 & Turbulent covariance (Reynolds stress) \\
$\mathbf{c}_{\text{auto}}$ & 3 & Auto-correlation peak center \\
$\mathbf{c}_{\text{cross}}$ & 3 & Cross-correlation peak center \\
\midrule
Total & 22 & \\
\bottomrule
\end{tabular}
\caption{Parameters of the stacked 3D Gaussian model.}
\end{table}

The Gaussian model for each correlation is:
\begin{equation}
    G(\mathbf{r}; A, B, \mathbf{c}, \bm{\Sigma}) = A \exp\left(-\frac{1}{2}(\mathbf{r} - \mathbf{c})^T \bm{\Sigma}^{-1} (\mathbf{r} - \mathbf{c})\right) + B
\end{equation}

The coupled constraint is enforced:
\begin{align}
    R_{\text{auto}}^{\text{model}}(\mathbf{r}) &= G(\mathbf{r}; A_{\text{auto}}, B_{\text{auto}}, \mathbf{c}_{\text{auto}}, \bm{\Sigma}_{\text{geo}}) \\
    R_{\text{cross}}^{\text{model}}(\mathbf{r}) &= G(\mathbf{r}; A_{\text{cross}}, B_{\text{cross}}, \mathbf{c}_{\text{cross}}, \bm{\Sigma}_{\text{geo}} + \bm{\Sigma}_{\text{turb}})
\end{align}

\subsection{Initial Guess via Moment Estimation}

A good initial guess is critical for nonlinear optimization. We estimate covariances from intensity-weighted second moments:
\begin{equation}
    \hat{\Sigma}_{ij} = \frac{\sum_{\mathbf{r}} w(\mathbf{r}) (r_i - \bar{r}_i)(r_j - \bar{r}_j)}{\sum_{\mathbf{r}} w(\mathbf{r})}
\end{equation}
where weights $w(\mathbf{r}) = \max(R(\mathbf{r}) - \tau, 0)$ focus on the peak region above threshold $\tau$.

The initial turbulent covariance is estimated as:
\begin{equation}
    \bm{\Sigma}_{\text{turb}}^{(0)} = \hat{\bm{\Sigma}}_{\text{cross}} - \hat{\bm{\Sigma}}_{\text{auto}}
\end{equation}
with eigenvalue clamping to ensure positive semi-definiteness.

\subsection{Weighted Least-Squares Optimization}

We minimize the weighted sum of squared residuals:
\begin{equation}
    \chi^2 = \sum_{\mathbf{r}} w_{\text{auto}}(\mathbf{r}) \left[R_{\text{auto}}(\mathbf{r}) - R_{\text{auto}}^{\text{model}}(\mathbf{r})\right]^2 + \sum_{\mathbf{r}} w_{\text{cross}}(\mathbf{r}) \left[R_{\text{cross}}(\mathbf{r}) - R_{\text{cross}}^{\text{model}}(\mathbf{r})\right]^2
\end{equation}

Weights are computed from normalized intensities:
\begin{equation}
    w(\mathbf{r}) = \sqrt{\frac{R(\mathbf{r})}{R_{\max}} + 0.1}
\end{equation}
giving higher importance to voxels near the peak where signal-to-noise is highest.

Optimization is performed using the trust-region reflective algorithm with bounds ensuring:
\begin{itemize}
    \item Amplitudes $A > 0$
    \item Diagonal covariance elements $\Sigma_{ii} > 0.1$ (preventing singularities)
    \item Peak centers within the ROI
\end{itemize}

\section{Results}

\subsection{Experimental Configuration}

We validate the method using synthetic data with:
\begin{itemize}
    \item Volume size: $64 \times 64 \times 16$ voxels
    \item Number of particles: 50 per frame
    \item Number of image pairs: 2000
    \item Camera angles: $\pm 45°$
\end{itemize}

The ground-truth Reynolds stress tensor (in voxels$^2$/frame$^2$):
\begin{equation}
    \mathbf{R}_{\text{true}} = \begin{pmatrix}
        3.0 & 0.9 & 0.3 \\
        0.9 & 2.4 & 0.6 \\
        0.3 & 0.6 & 1.5
    \end{pmatrix}
\end{equation}

\subsection{Reynolds Stress Recovery}

Table~\ref{tab:results} compares the fitted Reynolds stress components to ground truth.

\begin{table}[h]
\centering
\begin{tabular}{@{}lccccc@{}}
\toprule
Component & True & Fitted & Abs.\ Error & Rel.\ Error \\
\midrule
$R_{uu}$ & 3.00 & 3.96 & 0.96 & 32\% \\
$R_{vv}$ & 2.40 & 1.72 & 0.68 & 28\% \\
$R_{ww}$ & 1.50 & 2.67 & 1.17 & 78\% \\
$R_{uv}$ & 0.90 & 0.68 & 0.22 & 25\% \\
$R_{uw}$ & 0.30 & 0.36 & 0.06 & 19\% \\
$R_{vw}$ & 0.60 & 0.38 & 0.22 & 37\% \\
\midrule
\multicolumn{2}{l}{RMS Error:} & \multicolumn{2}{c}{0.57} & \\
\multicolumn{2}{l}{Mean Rel.\ Error:} & \multicolumn{2}{c}{33\%} & \\
\bottomrule
\end{tabular}
\caption{Reynolds stress recovery results for 2000 image pairs.}
\label{tab:results}
\end{table}

\subsection{Discussion}

Key observations:
\begin{enumerate}
    \item \textbf{Off-diagonal terms} ($R_{uv}, R_{uw}, R_{vw}$) are recovered with 19--37\% error, demonstrating sensitivity to shear stresses.

    \item \textbf{In-plane components} ($R_{uu}, R_{vv}$) show 28--32\% error, reasonable given the reconstruction and fitting challenges.

    \item \textbf{Out-of-plane component} ($R_{ww}$) has the largest error (78\%), attributable to the anisotropic volume resolution ($64 \times 64 \times 16$). With only 16 voxels in $z$, the correlation peak width in this direction is poorly resolved.

    \item \textbf{Mean displacement} is recovered as essentially zero ($< 0.01$ voxels), as expected for our zero-mean-flow test case.
\end{enumerate}

\section{Conclusions}

We have presented a complete pipeline for extracting Reynolds stress tensors from stereo PIV data using ensemble correlation analysis. The method:

\begin{enumerate}
    \item Generates synthetic stereo images with prescribed Reynolds stress statistics
    \item Reconstructs 3D volumes using MinLOS to suppress ghost artifacts
    \item Accumulates ensemble correlations with background subtraction
    \item Fits a coupled 22-parameter Gaussian model to extract geometric and turbulent covariances
\end{enumerate}

Validation with 2000 synthetic image pairs demonstrates Reynolds stress recovery with approximately 30\% mean relative error. The primary limitation is anisotropic volume resolution, particularly affecting the out-of-plane stress component $R_{ww}$.

Future improvements could include:
\begin{itemize}
    \item Increased z-resolution through additional camera views
    \item More sophisticated ghost suppression algorithms
    \item Regularization or prior constraints on the Reynolds stress tensor
    \item Extension to spatially-varying stress fields
\end{itemize}

\appendix

\section{Algorithm Summary}

\begin{algorithm}[H]
\caption{Stereo Ensemble PIV for Reynolds Stress Extraction}
\begin{algorithmic}[1]
\Require Stereo image pairs $\{(I_1^A, I_2^A, I_1^B, I_2^B)_k\}_{k=1}^N$
\Ensure Reynolds stress tensor $\mathbf{R}$
\State Initialize FFT accumulators: $\hat{S}_{AA} \gets 0$, $\hat{S}_{AB} \gets 0$
\State Initialize volume accumulators: $S_A \gets 0$, $S_B \gets 0$
\For{$k = 1$ to $N$}
    \State $V_A \gets \text{MinLOS}(I_1^A, I_2^A)$ \Comment{Reconstruct frame A}
    \State $V_B \gets \text{MinLOS}(I_1^B, I_2^B)$ \Comment{Reconstruct frame B}
    \State $\hat{S}_{AA} \gets \hat{S}_{AA} + |\mathcal{F}(V_A)|^2$
    \State $\hat{S}_{AB} \gets \hat{S}_{AB} + \mathcal{F}(V_A) \cdot \mathcal{F}(V_B)^*$
    \State $S_A \gets S_A + V_A$, $S_B \gets S_B + V_B$
\EndFor
\State $\langle R_{AA} \rangle \gets \mathcal{F}^{-1}[\hat{S}_{AA}/N]$ \Comment{Ensemble correlations}
\State $\langle R_{AB} \rangle \gets \mathcal{F}^{-1}[\hat{S}_{AB}/N]$
\State $R_{AA}^{\text{bg}} \gets (S_A/N) \star (S_A/N)$ \Comment{Background}
\State $R_{AB}^{\text{bg}} \gets (S_A/N) \star (S_B/N)$
\State $\tilde{R}_{AA} \gets \langle R_{AA} \rangle - R_{AA}^{\text{bg}}$ \Comment{Subtract background}
\State $\tilde{R}_{AB} \gets \langle R_{AB} \rangle - R_{AB}^{\text{bg}}$
\State $(\bm{\Sigma}_{\text{geo}}, \bm{\Sigma}_{\text{turb}}) \gets \text{FitStackedGaussian}(\tilde{R}_{AA}, \tilde{R}_{AB})$
\State \Return $\mathbf{R} \gets \bm{\Sigma}_{\text{turb}}$
\end{algorithmic}
\end{algorithm}

\end{document}
